% ====================================================================
\chapter{How to Use This Book}
% ====================================================================

This book is a self-contained textbook for the opening part of ECN115
\emph{Mathematical Methods in Economics and Finance}. It covers everything presented in the first two
lectures of the module, which together make up Topic~1 (\emph{Introduction}) and the first steps into
Topic~2 (\emph{Functions}) and Topic~3 (\emph{Limits and smooth functions}). It is written so that a
student who has never seen the lecture slides can learn the material from these pages alone, and so that
a student who has seen them can use it to understand every step properly rather than memorise it.

The slides present ideas in the order in which they were spoken. A textbook can do better: here the
material has been reorganised so that every idea is introduced only after the ideas it depends on. For
example, quantifiers (``for all'' and ``there exists'') appear in the second lecture but are needed
much earlier to state definitions precisely, so they are taught in \cref{ch:logic} alongside the rest
of the language of mathematics.

\section*{How each chapter is organised}
Every chapter follows the same pattern, so that you always know where to find things.
\begin{enumerate}[label=\arabic*.]
  \item \textbf{Introduction and prerequisites.} A short paragraph saying what the chapter is about and
  why it matters, followed by a grey box listing what you need to know first, with references back to
  earlier chapters.
  \item \textbf{The core material.} Definitions, explanations, results and proofs, developed in prose.
  Each important idea is explained four ways: what it says (the definition), what it means (intuition),
  why it matters (its role in the module and in economics), and how to use it (worked examples).
  \item \textbf{Worked examples.} Every example from the lectures is reproduced and explained in full,
  usually with more steps shown than on the slides. Further examples have been added wherever an idea
  needed more practice.
  \item \textbf{Chapter review.} A summary in prose, the key definitions and key equations collected
  in one place, connections to other chapters, the common mistakes, and an exam checklist.
  \item \textbf{Practice questions with full solutions.} Questions are labelled by type:
  \emph{basic} (direct use of a definition), \emph{conceptual} (explain why), \emph{mathematical}
  (calculation or proof), \emph{application} (an economic or applied setting), \emph{multi-step}
  (several ideas combined) and \emph{exam-style} (written in the format of the past ECN115 papers).
  Every question uses only material covered up to that point. Full solutions follow the questions.
\end{enumerate}

\section*{The boxes}
Different kinds of content are set in different boxes so that you can see at a glance what you are
reading.

\medskip
\begin{tabularx}{\textwidth}{@{}p{46mm}X@{}}
\boxkey{defcol}{defbg}{Definition} & A precise statement of what a term means. Definitions must be learned
exactly: past ECN115 papers award marks for stating them.\\[1.2ex]
\boxkey{thmcol}{thmbg}{Theorem / Proposition} & A statement that has been, or will be, proved from the
definitions and axioms.\\[1.2ex]
\boxkey{excol}{exbg}{Worked Example} & A problem solved step by step, with the reasoning behind each step.\\[1.2ex]
\boxkey{corecol}{corebg}{CORE CONCEPT} & An idea on which a large part of the module rests. If you are
short of time, make sure you understand every box with this label.\\[1.2ex]
\boxkey{formcol}{formbg}{Key Formula} & An important formula, given the full treatment: what each symbol
means, where it comes from, when to use it, an example, common mistakes, rearrangements, what happens when
the inputs change, and how it appears in questions.\\[1.2ex]
\boxkey{warncol}{warnbg}{Common mistake} & An error that students frequently make, and how to avoid it.\\[1.2ex]
\boxkey{examcol}{exambg}{Recognising this in an exam} & How a type of question is usually worded and what
it is really asking.\\[1.2ex]
\boxkey{flagcol}{flagbg}{Note on the source material} & A place where the slides or the study guide
contain an error, a typographical slip, or an ambiguity. The correct version is given, and nothing is
silently changed.\\[1.2ex]
\boxkey{extcol}{extbg}{Beyond the course materials} & Knowledge that is not in the ECN115 materials but
helps you understand them. You are not expected to reproduce it unless your lecturer says otherwise.\\[1.2ex]
\boxkey{notecol}{notebg}{Note} & A remark, a qualification or a piece of context.\\
\end{tabularx}

\section*{Source tags}
Small grey tags such as \srcin{L1, slide 17} at the right-hand edge of a box tell you where the material
comes from: \emph{L1} and \emph{L2} are the two lecture slide sets, \emph{SG} is the Topic~1 study and
revision guide, and \emph{Jan 2024} and similar labels are past examination papers. Use these tags to
find the original when you want to compare.

\section*{A note on examinability}
This book never claims that a particular item will be examined. Where the course materials themselves
say something about assessment it is reported: for example, the slides describe the Least Upper Bound
property as ``extra material, non-examinable''. Where past papers show that a type of question has been
set before, the book says so, because that is a matter of record; it is not a prediction.

\section*{Suggested way to study}
On a first reading, work through each chapter in order and attempt every worked example yourself before
reading its solution. On a second reading, concentrate on the CORE CONCEPT boxes, the key formulas and
the chapter reviews, and attempt the practice questions with the solutions covered. Before an assessment,
use the \emph{Final Master Review} at the end of the book: it contains a summary of the whole module, a
concept map, a master formula sheet, a glossary, a checklist of everything you should be able to do, and
a final set of mixed practice questions with solutions.

% ====================================================================
\chapter{Module Overview}
% ====================================================================

\section*{The module}
ECN115 \emph{Mathematical Methods in Economics and Finance} is taught at Queen Mary University of London in
2026--27 by Dr Evgenii Safonov (e-mail \texttt{e.safonov@qmul.ac.uk}, office GC332, office hours Tuesdays
15:00--17:00). Lectures take place on Wednesdays 09:00--11:00 (Arts Two lecture theatre) and Thursdays
09:00--11:00 (Bancroft Mason lecture theatre); classes are listed on MyTimetable, and undergraduate support
is through AskQM. \src{L1, slides 2--6}

Assessment consists of a one-hour closed-book midterm test worth 30\% (on 11 November 2026) and a two-hour
closed-book final examination worth 70\% (in January). Calculators were not permitted on any of the three
past papers summarised in the study guide. Weekly homework is not graded but is described as an important
component: each week you attempt the homework, the most challenging problems are analysed in class, and
then solutions are posted. As the slides put it, ``it is almost impossible to learn mathematics without
practice''.

No textbook is required. The slides suggest three for supplementary reading: Sydsæter, Hammond and others,
\emph{Essential Mathematics for Economic Analysis}; Simon and Blume, \emph{Mathematics for Economists}; and,
for more advanced study, Sydsæter, Hammond, Seierstad and Strøm, \emph{Further Mathematics for Economic
Analysis}.

\section*{The six topics of the module}
\begin{enumerate}[label=\textbf{Topic \arabic*.}, leftmargin=5.5em]
  \item \textbf{Introduction.} Mathematical arguments. Sets, numbers, vectors and sequences. Absolute value,
  inequalities, mathematical induction, the binomial formula.
  \item \textbf{Functions.} The concept of a function. Elementary numerical functions. Basic properties of
  functions. Relationships between functions.
  \item \textbf{Limits and smooth functions.} Limits of sequences. Continuity and differentiability of
  functions. Properties of the derivative. The chain rule and other rules of differentiation. Higher-order
  derivatives.
  \item \textbf{Univariate optimisation.} The existence of maxima and minima. The first-order condition for
  optimality. Corner solutions. Dependence of the solution on the parameter.
  \item \textbf{Linear algebra.} Systems of linear equations. Matrix algebra. Matrix inversion.
  \item \textbf{Multivariate analysis and optimisation.} Multivariate functions. Partial derivatives.
  Unconstrained and constrained optimisation: existence, first-order conditions.
\end{enumerate}

\section*{What this book covers}
The materials provided for this book are the slides of Lectures 1 and 2, a study and revision guide for
Topic~1 (which also summarises the January 2024, 2025 and 2026 papers), and nothing else. The book therefore
covers exactly the content of those two lectures: all of Topic~1 except vectors, the concept of a function
and polynomials from Topic~2, and the definition of the limit of a sequence from Topic~3. The slides list
the plan for the following two lectures (properties of limits, calculating limits, infinite sums and an
introduction to derivatives in Lecture~3; more functions and derivatives in Lecture~4), but no material
from those lectures was provided, so it is not covered here.

\begin{sourcenote}[Coverage]
The Topic~1 syllabus lists \emph{vectors}, but vectors do not appear in the slides of Lectures 1 and 2.
They are therefore not taught in this book. If vectors are covered later in the module, treat that material
as an addition to \cref{ch:sets,ch:numbers}.
\end{sourcenote}

\section*{How the chapters fit together}
The chapters build on one another as shown in the map below. The first three chapters provide the
language (logic and sets) and the objects (real numbers and their properties). Chapters \ref{ch:ineq} and
\ref{ch:abs} use the properties of real numbers to solve inequalities, which is the single most frequently
used skill in the module. Chapters \ref{ch:sum}--\ref{ch:binom} develop sums and proof by induction, and
use them to prove the binomial formula. Chapter~\ref{ch:func} defines functions and proves that a polynomial
can be factorised through each of its roots. Chapter~\ref{ch:limits} brings everything together in the
definition of the limit of a sequence, which uses quantifiers, absolute values, inequalities and the
triangle inequality all at once, and which is the foundation of Topic~3.

\begin{figure}[ht]
\centering
\begin{tikzpicture}[node distance=9mm and 6mm, every node/.style={mapnode=34mm}]
  \node (logic) {1. Language of\\ mathematics};
  \node[right=of logic] (sets) {2. Sets};
  \node[right=of sets] (num) {3. Real numbers\\ and their axioms};
  \node[below=of num] (ineq) {4. Inequalities};
  \node[left=of ineq] (abs) {5. Absolute value,\\ triangle inequality};
  \node[left=of abs] (sum) {6. Summation\\ notation};
  \node[below=of sum] (ind) {7. Mathematical\\ induction};
  \node[right=of ind] (bin) {8. Binomial\\ formula};
  \node[right=of bin] (fun) {9. Functions and\\ polynomials};
  \node[below=of bin, mapnode=60mm, fill=corebg, draw=corecol] (lim) {10. Sequences and limits};
  \draw[maparrow] (logic) -- (sets);
  \draw[maparrow] (sets) -- (num);
  \draw[maparrow] (num) -- (ineq);
  \draw[maparrow] (ineq) -- (abs);
  \draw[maparrow] (num) to[bend right=12] (sum);
  \draw[maparrow] (sum) -- (ind);
  \draw[maparrow] (ind) -- (bin);
  \draw[maparrow] (ind) to[bend right=25] (fun);
  \draw[maparrow] (abs) -- (lim);
  \draw[maparrow] (fun) -- (lim);
  \draw[maparrow] (logic) to[out=-120,in=150] (ind);
\end{tikzpicture}
\caption*{\textbf{Map of the book.} How the chapters depend on one another. An arrow from one chapter to another means that the
second uses ideas from the first. The limit of a sequence (\cref{ch:limits}) draws on almost everything
before it.}
\end{figure}

\section*{Where the lectures appear in this book}
\begin{center}
\begin{tabularx}{\textwidth}{@{}Xlc@{}}
\toprule
Lecture material & Slides & Chapter\\
\midrule
Elementary logic: propositions, proof, axiom, theorem, $\Rightarrow$, $\Leftrightarrow$, negation & L1, 8 & \ref{ch:logic}\\
``There exists'' and ``for all'' & L2, 21 & \ref{ch:logic}\\
Sets, subsets, intersection, union, set-builder notation, set difference & L1, 10--14 & \ref{ch:sets}\\
Number systems, properties O1--O2, A1--A5, M1--M6, Least Upper Bound, intervals & L1, 16--20 & \ref{ch:numbers}\\
Basic and quadratic inequalities (Examples 1--3) & L1, 22--25 & \ref{ch:ineq}\\
Absolute value, triangle inequality & L2, 10--11 & \ref{ch:abs}\\
Summation operator and its properties & L1, 27--30 & \ref{ch:sum}\\
Mathematical induction & L1, 32--36; L2, 4--8 & \ref{ch:ind}\\
The binomial formula, factorials, binomial coefficients & L2, 13--14 & \ref{ch:binom}\\
Functions, polynomials, Theorem L2.1 & L2, 16--19 & \ref{ch:func}\\
Sequences, numerical sequences, limits & L2, 22--30 & \ref{ch:limits}\\
\bottomrule
\end{tabularx}
\end{center}

\section*{Why mathematics in an economics degree?}
Economics makes claims about how people, firms and governments behave and about the consequences of
that behaviour. Many of these claims are precise enough to be written down mathematically: a consumer
chooses the bundle that maximises utility subject to a budget, a firm chooses output to maximise profit,
an interest rate compounds over time, a sequence of quarterly output figures approaches a long-run level.
Writing such claims mathematically has two great advantages. First, it forces precision: a statement such
as ``demand falls when the price rises'' becomes an inequality that can be checked. Second, it allows
\emph{proof}: once the assumptions are written down, their consequences can be derived with certainty
rather than argued about. The material in this book is the toolkit for both. Inequalities describe
constraints and comparisons; sets describe the options available to a decision-maker; sums describe
totals over time or over people; induction proves statements about every period at once; functions
describe relationships such as cost as a function of output; and limits describe long-run behaviour.
